\documentclass[showpacs,amsmath,amssymb,twocolumn,superscriptaddress,notitlepage,preprintnumbers,aps,pra,10pt,floatfix]{revtex4-2}
\usepackage{graphicx}% Include figure files
\usepackage{dcolumn}% Align table columns on decimal point
\usepackage{bm}% bold math
\usepackage{qcircuit}
\usepackage{braket}
\usepackage{color}
\usepackage{amsthm}
\usepackage{chemarrow}
\usepackage{overpic}
\usepackage{diagbox}
\usepackage{braket}
\usepackage{subfigure}
\usepackage{makecell}
\newtheorem{theorem}{Theorem}
\newtheorem{definition}{Definition}
\newtheorem{lemma}{Lemma}
\newtheorem{corollary}{Corollary}
\usepackage[colorlinks=true,linkcolor=red,bookmarks=true,breaklinks=true]{hyperref}
\usepackage{chngcntr}
% \newcounter{pretheorem}
% \counterwithin{theorem}{pretheorem}
% \renewcommand\thetheorem{\arabic{pretheorem}\alph{theorem}}
% \newcommand{\theoremgroup}{\refstepcounter{pretheorem}}
\begin{document}% Start the document part
\let\oldacl\addcontentsline
\renewcommand{\addcontentsline}[3]{}
\newcommand{\Tr}{\mathrm{Tr}}
\newcommand{\szx}[1]{\textcolor{red}{(SZX: #1)}}
\title{What enables efficient classical simulation of noisy quantum circuits}
% \author{Zhong-Xia Shang}
% % \email[]{ustcszx@mail.ustc.edu.cn}
% % % \affiliation{QICI Quantum Information and Computation Initiative, Department of Computer Science,
% % % The University of Hong Kong, Pokfulam Road, Hong Kong}
% % \affiliation{HK Institute of Quantum Science $\&$ Technology, The University of Hong Kong, Hong Kong, China}
% % \affiliation{QICI Quantum Information and Computation Initiative, Department of Computer Science,
% % The University of Hong Kong, Hong Kong, China}
% % \affiliation{Shanghai Research Center for Quantum Science and CAS Center for Excellence in Quantum Information and Quantum Physics,
% % University of Science and Technology of China, Shanghai 201315, China}
% \author{Si-Yuan Shen...}

\begin{abstract}
The motivation is to understand what fundamental properties of noisy quantum circuits make efficient classical simulation possible.
\end{abstract}
\maketitle

% \setcounter{tocdepth}{1}
% \tableofcontents
% \noindent\textit{\textbf{Inverse polynomial scaling of Pauli spectra.—}}
\section{Inverse polynomial scaling of Pauli spectra}

Given a noisy quantum circuit $\mathcal{C}[\cdot]$, the goal is to estimate the expectation value:
\begin{eqnarray}
\Tr(\mathcal{C}[|0\rangle\langle 0|^{\otimes n}]O)=\Tr(|0\rangle\langle 0|^{\otimes n}\mathcal{C}^\dag[O])
\end{eqnarray}
The idea of Pauli path simulation is that $\mathcal{C}^\dag[O]$ has the Pauli decomposition:
\begin{eqnarray}
O_H=\mathcal{C}^\dag[O]=\sum c_i P_i
\end{eqnarray}
satisfying $|c_i|\geq |c_j|$ if $i\leq j$, and we hope to use another operator $\widetilde{O_H}$ with the form:
\begin{eqnarray}
\widetilde{O_H}=\mathcal{C}^\dag[O]=\sum_{i=0}^{k-1} b_i P_i
\end{eqnarray}
to approximate $O_H$.

Under various randomness assumptions such as fixed circuit with random input, the average error of using $\widetilde{O_H}$ to replace $O_H$ for expectation value estimation can be related to Frobenius norm:
\begin{eqnarray}
E[\Tr(|0\rangle\langle 0|^{\otimes n} (O_H-\widetilde{O_H}))^2]=x \|\Delta O_H\|_F^2=\epsilon^2
\end{eqnarray}
with $\Delta O_H=O_H-\widetilde{O_H}$.

Efficient classical simulation requires the classical complexity $C$ has the scaling:
\begin{eqnarray}
C=\mathcal{O}(\text{poly}(\epsilon^{-1}))
\end{eqnarray}
Since $C\geq k$, as a necessary condition, we require:
\begin{eqnarray}
k=\mathcal{O}(\text{poly}(\epsilon^{-1}))
\end{eqnarray}
For fixed $k$, the best we can expect or the smallest $\|\Delta O_H\|_F$ we can have is:
\begin{eqnarray}
f(k)=\min \|\Delta O_H\|_F^2=\sum_{i=k}c_i^2
\end{eqnarray}
Suppose $k= y \varepsilon^{-r}+o(\varepsilon^{-r})$, we have:
\begin{eqnarray}
f(k)=\sum_{i=k}c_i^2=\frac{\epsilon^2}{x}= \frac{1}{x}(\frac{y}{k})^{2/r} + o(k^{-2/r})
\end{eqnarray}
whose derivative therefore has the relation:
\begin{eqnarray}
-\frac{d f}{dk}=c_k^2=\frac{2y^{2/r}}{rx}k^{-1-2/r} + o(k^{-1-2/r})
\end{eqnarray}
Since $r\geq0$, this leads to:
\begin{eqnarray}
c_k^2=\mathit{o}(k^{-1})
\end{eqnarray}
Therefore, a necessary condition for efficient classical simulation is to have $O_H$ with an inverse polynomial Pauli spectrum. The only assumption in the derivations is the randomness that enables the reduction from the error of estimation to the Frobenius norm of $\Delta O_H$. (see Fig. \ref{fig:necessary_condition}.)

\begin{figure}[htbp!]
    \centering
    \includegraphics[width=1\linewidth]{pic/necessary_condition.png}
    \caption{Noiseless circuit decay slowly compare to noisy circuit. Effective truncation method exist only if the decay of $c_i^2$ is faster than $i^{-1}$}
    \label{fig:necessary_condition}
\end{figure}

\section{Remained problems and what to do next?}

\subsection{Use entropy to reformulate the result}
We can normalize the spectra $\{c_i^p\}$ to a probability distribution. Then the inverse polynomial scaling can also imply some properties of entropies in terms of this probability distribution.

\subsection{Various randomness assumption}
Reduction from the error of estimation to the Frobenius norm of $\Delta O_H$ is due to randomness. There are various randomness assumptions: random circuit, random sampling directions, random input, or even random operators? For the random circuit case, how to formalize our results? The average Pauli spectra is an inverse polynomial?

Also, what kind of randomness can lead to Frobenius norm (a dictionary)? However, even if a randomness leads to other norms of $\Delta O_H$, the derivations should also apply.

\subsection{Towards fixed circuit}
Consider fixed circuit and fixed input. Then:
\begin{eqnarray}
\langle O\rangle =\sum c_i \Tr(P_i|0\rangle\langle 0|^{\otimes n})=\sum_i v_i
\end{eqnarray}
Can we say that $v_i$ should have inverse polynomial scaling to enable efficient classical simulation? The problem here is that $v_i$ is not positive and we have no reason to use $v_i^2$ as in the random case with Frobenius norm. 

Another issue in fixed circuits is that the input itself can preclude many Pauli operators: $|0\rangle$ preclude $X$ and $Y$.

\subsection{Other rules}
Even if the spectra is not inverse polynomial, the problem still has the chance to be efficiently simulated as long as $\|\Delta O_H\|_F^2$ is already exponentially small. In other word, even if $k=\mathcal{O}(\text{poly}(\epsilon^{-1}))$ is not satisfied, $\epsilon=\mathcal{O}(\exp(-n))$ is enough. This corresponds to the case of the output state closing to the identity operator?

\subsection{Gap between our result and practical algorithms}
The spectra result is only a necessary condition and it only relates the spectra of the output operator, so it is an operator-related property rather than a circuit-related. However, in practical Pauli classical algorithms, we consider Pauli path, where the depth and the intermediate Paulis are also important. Is it possible to introduce a time register to also account for the influence of depth?

\subsection{Possible numerical simulation}
We can do various numerical simulations to see whether there is a complexity phase transition (non-inverse poly to inverse poly) when gradually adding noise to a circuit. Setting can be various: brick random circuit, Clifford+T circuit...
Maybe the spectra can also reveal chaotic behaviors: GUE to Poisson distribution.

\subsection{operator 1 norm is important?}
Does the 1 norm $\sum_i |c_i|$ also play an important role in classical simulation? The observation is that the Clifford circuit won't change 1 norm, the T gate can increases the 1 norm to even exponentially large, and the depolarizing noise can reduce the 1 norm to even exponentially small. Therefore the value of 1 norm characterize the battle between magic and noise. Thus, is the classical complexity of simulation algorithms also related to this value? What's the relation between 1 norm and the inverse polynomial spectra?

\subsection{Truncation and sampling separation?}
Consider a noisy circuit: all 0 input state, amplitude sampling noise, and all $Z$ operators. In this case, $O_H$ has $2^n$ Pauli operators containing only $I$ and $Z$ with all $c_i=2^{-n}$. The spectra can not be bounded by $i^{-1}$ and therefore the Pauli truncation-based methods require $k$ to be exponential to have an accurate approximation. However, we can also use the randomized method to sample Pauli operators, in this case, since each sampled Pauli operator has the expectation value equal to 1, the exact result. Therefore, there is no variance, and a small number of samples is already enough. For sampling-based methods, it seems that flat spectra correspond to small variance, even though this spectra is hard for truncation methods. What's the origin of this difference and can be unify them in a single picture (measure)? 



\subsection{TBA}
\bibliography{ref}
\clearpage
\begin{appendix}
\onecolumngrid
\renewcommand{\addcontentsline}{\oldacl}
\renewcommand{\tocname}{Appendix Contents}
\renewcommand\thetheorem{\arabic{pretheorem}}
\tableofcontents

\section{Preliminaries}
\subsection{Vectorization}
\begin{definition}[Vectorization]
Given a $n$-qubit matrix $O=\sum_{ij}o_{ij}|i\rangle\langle j|$, the vectorization mapping $\mathcal{V}$ has the operation: $\mathcal{V}[O]=\ket{O\rangle}=\sum_{ij}o_{ij}|i\rangle|j\rangle$. 
\end{definition}
Under the definition of vectorization, we have
\begin{eqnarray}
\Tr(B^\dag A)=\braket{\langle B|A\rangle}.
\end{eqnarray}
Therefore, Pauli operators will become orthogonal vectors under the vectorization mapping. Indeed, we have the Pauli-Bell correspondence
\begin{eqnarray}
I=\begin{pmatrix}1& 0\\0&1\end{pmatrix}&&\xrightarrow[]{\mathcal{V}}\sqrt{2}|\Phi^+\rangle=|0\rangle|0\rangle+|1\rangle|1\rangle\nonumber,\\
Z=\begin{pmatrix}1& 0\\0&-1\end{pmatrix}&&\xrightarrow[]{\mathcal{V}}\sqrt{2}|\Phi^-\rangle=|0\rangle|0\rangle-|1\rangle|1\rangle\nonumber,\\X=\begin{pmatrix}0& 1\\1&0\end{pmatrix}
&&\xrightarrow[]{\mathcal{V}}\sqrt{2}|\Psi^+\rangle=|0\rangle|1\rangle+|1\rangle|0\rangle\nonumber,\\
Y=\begin{pmatrix}0& -i\\i&0\end{pmatrix}&&\xrightarrow[]{\mathcal{V}}-i\sqrt{2}|\Psi^-\rangle=|0\rangle|1\rangle-|1\rangle|0\rangle.
\end{eqnarray}

\subsection{Various norms}






\section{Heisenberg Pauli truncation}
Given a $n$-qubit noisy quantum circuit manifested as a CPTP map: $\mathcal{C}[\cdot]$ acting on an initial state $|\psi\rangle$, the goal is to estimate the expectation value of an operator $O$
\begin{eqnarray}
\langle O\rangle=\Tr(\mathcal{C}[|\psi\rangle\langle \psi|]O)=\Tr(|\psi\rangle\langle \psi|\mathcal{C}^\dag[O]).
\end{eqnarray}
We define $O_H=\mathcal{C}^\dag[O]$ and decompose it in the Pauli basis
\begin{eqnarray}
O_H=\mathcal{C}^\dag[O]=\sum_{i=1}^M c_i P_i,
\end{eqnarray}
where $\{c_i\}$ are arranged in descending order in terms of the absolute values. 
Since $M$ can generally be exponentially large in qubit number $n$, the idea of truncation-based methods is to use another operator $\widetilde{O_H}$ to approximate $O_H$.  $\widetilde{O_H}$  has the form
\begin{eqnarray}
\widetilde{O_H}=\mathcal{C}^\dag[O]=\sum_{j=1}^{K-1} b_j P_j,
\end{eqnarray}
where $K = \text{poly}(n)$ and each term $b_j$ can be calculated within polynomial time in $n$.

Under the vectorization picture, we have
\begin{eqnarray}\label{b3}
\langle O\rangle=(\langle \psi|\otimes \langle \psi^*|)\ket{O_H\rangle},
\end{eqnarray}
with $\ket{O_H\rangle}=\mathcal{V}[O_H]$ normalized. We also define normalized $|\widetilde{O_H}\rangle\!\rangle=\mathcal{V}[\widetilde{O_H}]$.

Consider an input state ensemble $\mu$ with $\|\mathbb{E}_{\psi\in\mu}|\psi\rangle\langle \psi|\|_{\infty}=2^{-n}g$ and define $\Delta O_{H}=O_H-\widetilde{O_H}$ and ${|\Delta O_{H}\rangle\!\rangle}=|{O_H\rangle\!\rangle}-|{\widetilde{O_H}\rangle\!\rangle}$, we have the average approximation error
\begin{eqnarray}\label{eq:hpt}
\mathbb{E}_{\psi\in\mu}\Tr(|\psi\rangle\langle\psi| \Delta O_H)^2&&=\mathbb{E}_{\psi\in\mu}(\langle\psi|\otimes \langle\psi^*| )|{\Delta O_{H}\rangle\!\rangle}{\langle\!\langle\Delta O_{H}|}(|\psi\rangle\otimes |\psi^*\rangle)\nonumber\\
\leq && \mathbb{E}_{\psi\in\mu}\Tr((|\psi\rangle\langle \psi| \otimes I_n)|{\Delta O_{H}\rangle\!\rangle}{\langle\!\langle\Delta O_{H}|})\nonumber\\
\leq && \|\mathbb{E}_{\psi\in\mu}|\psi\rangle\langle \psi|\|_{\infty}{\langle\!\langle\Delta O_{H}|\Delta O_{H}\rangle\!\rangle}\nonumber\\
=&&g \|{\Delta O_{H}\rangle\!\rangle}\|_{n2}^2,
\end{eqnarray}
where the normalized vector-2 norm is defined as $\|\cdot\|_{n2}=2^{-n/2}\|\cdot\|_{2}$. For example, $\|\ket{O_H\rangle}\|_{n2}=\sum_i c_i^2$. When $\mu$ is a state 1-design, we have $g=1$ and the average approximation error is directly related to $\|\ket{\Delta O_{H}\rangle}\|_{n2}^2$.

%%%%%%%%%%%%%%%%%%%%%%%%%%%%%%%%%%%%%%%%%%
\section{Pauli spectra phase transition diagnoses classical simulation ability}

Ignoring the computational complexity that realizes such a truncation, we can regards $K$ as the com complexity of Pauli truncation-based classical simulation methods. 
In this case, the best truncation strategy we can expect, which has the smallest approximation error, is making all $b_i$ equal to $c_i$ and keeping $K-1$ largest Pauli terms in $\overline{O_H}$. This gives truncated operator $\Delta O_H$. 
According to Eq.~\eqref{eq:hpt}, this efficient classical simulation achieves an error on average approximation such that
\begin{eqnarray}
\text{err} =\mathbb{E}_{\psi\in\mu}\Tr(|\psi\rangle\!\langle\psi| \Delta O_H)^2 \leq  g f(K),
\end{eqnarray}
where we define
\begin{eqnarray}
\|\ket{\Delta O_{H}\rangle}\|_{n2}^2=\sum_{i=K}^M c_i^2:=f(K).
\end{eqnarray}
We also imply that $c_i^2 > c_{i+1}^2$ and $f(\cdot)$ is the function that map the  
To ensure that truncation algorithm is practical, we prove the necessary condition that there exist a truncation algorithm within $\textbf{FPTAS}$ complexity class.
\begin{lemma}
\label{lemma:FPTAS_condition}
   If there exist a truncation algorithm within $\textbf{FPTAS}$ class, then the function $f$ must satisfied $ f\left( \text{poly}(B)K \right) = \varepsilon^2 \cdot O(1/B^2)$ for $ f\left( K \right) = \varepsilon^2 $ and arbitrary $B>0$ that closed to $0$ .
\end{lemma}
\begin{proof}
    Given the existence of $\textbf{FPTAS}$ algorithm, there exists an algorithm with time complexity $\text{poly}(n) \cdot \text{poly}({1}/{C})$ for any inputs $C > 0$. This algorithm evaluates only $K-1 \sim \text{poly}(n)\cdot \text{poly}({1}/{C})$ numbers of non-zero $b_j$. Besides,  $\{b_j, P_j\}_{j =1}^{M}$ satisfy
    \begin{align}
     |(1\pm C)\langle O \rangle - \langle O \rangle|^2&=  \sum_{i=1}^M (c_i - b_j)^2 \\
     &\geq  \sum_{i = K}^{M} c_i^2 = f(\text{poly}(n) \cdot \text{poly}({1}/{C})).
    \end{align}
    Through taking $C = 1/B$ and $\langle O \rangle^2 \leq 1$, we conclude that $f(\text{poly}(n) \cdot \text{poly}~B) = O(1/B^2)$.
\end{proof}
Supposed the leading order in $\text{poly}(B)$ is $\eta$, we conclude that $ f\left( O(B^\eta) K \right) = \varepsilon^2 \cdot O(1/B^2)$. This requirement needs
\begin{eqnarray}
f(x)= O(x^{-\frac{2}{\eta}}).
\end{eqnarray}
In other words, efficient classical simulation indicates that $f(x)$ is bounded by a negative power function with certain constant, which we denote $A$.

\subsection{Ideal case}
Here, we assume $f(x)$ is exactly equal to $f(x)=Ax^{-\frac{2}{\eta}}$, then we have
\begin{eqnarray}
-\frac{df(x)}{dx}=c_x^2=\frac{2}{\eta}Ax^{-(1+\frac{2}{\eta})}=A'x^{-(1+\frac{2}{\eta})}=g(x),
\end{eqnarray}
with $A'=\frac{2}{\eta}A$. Therefore, efficient classical simulation also indicates the Pauli spectra $\{c_i^2\}$ is bounded by a negative smaller than -1 power function with a constant $A'$.

The phase transition to  $\textbf{FPTAS}$ complexity class can be diagnosed by numerical tests. Here, we show the strategies. 
\begin{theorem}
    If there exist a truncation algorithm within $\textbf{FPTAS}$ class, then the function $C(\alpha):=\text{log}~\max_{i = 1}^{M} \{c_i^2 i^{\alpha}\} $ has a first-order phase transition point $\alpha_0 >1$.  
\end{theorem}
\begin{proof}
    With Lemma~\ref{lemma:FPTAS_condition},  $f(x)=O(x^{-\frac{2}{\eta}})$, then we have
\begin{eqnarray}
-\frac{df(x)}{dx}=c_x^2=O(x^{-(1+\frac{2}{\eta})}).
\end{eqnarray}
Therefore, efficient classical simulation also indicates the Pauli spectra $\{c_i^2\}$ is bounded by a negative smaller than -1 power function with a constant $A'$. Then we conclude
\begin{align}
    C(\alpha)&:=\max_{i = \{1,\cdots,M\}} \{c_i^2 i^{\alpha}\} \\
    &=\max_{i = \{1,\cdots,M\}} O(i^{\alpha-(1+\frac{2}{\eta})}).
\end{align}
\begin{eqnarray}\label{pt1}
C(\alpha)= \begin{cases}
O(1)& \quad \alpha\leq \frac{2}{\eta} +1, \\
\text{Const} \times M^{\alpha}& \quad \alpha>\alpha_0>\frac{2}{\eta}+1,\\
\end{cases}
\end{eqnarray}
In other words, There exists a a first-order phase transition point $\alpha_0 > 1 + 2/\eta$  such that
\begin{eqnarray}
\frac{\partial\log(C(\alpha))}{\partial \alpha}= \begin{cases}
0 & \quad \alpha\leq \frac{2}{\eta}+1, \\
\log(M)& \quad \alpha>\alpha_0.\\
\end{cases}
\end{eqnarray}
\end{proof}


\szx{S version}We can try to use another function $h(\alpha, x)$ to bound $g(x)$
\begin{eqnarray}
h(\alpha, x)=C(\alpha)x^{-(1+\alpha)}\geq g(x),
\end{eqnarray}
leading to
\begin{eqnarray}
C(\alpha)=\max_x B x^{\alpha-\frac{2}{\eta}}=\max_x\frac{2}{\eta}A^{\frac{2}{\eta}}x^{\alpha-\frac{2}{\eta}}.
\end{eqnarray}
For different regimes, we have
\begin{eqnarray}\label{pt1}
C(\alpha)= \begin{cases}
\frac{2}{\eta}A^{\frac{2}{\eta}}x^{\alpha-\frac{2}{\eta}}\mid_{x=1}=\frac{2}{\eta}A^{\frac{2}{\eta}} & \quad \alpha\leq \frac{2}{\eta}, \\
\frac{2}{\eta}A^{\frac{2}{\eta}}x^{\alpha-\frac{2}{\eta}}\mid_{x=M}=\frac{2}{\eta}A^{\frac{2}{\eta}}M^{\alpha-\frac{2}{\eta}}& \quad \alpha>\frac{2}{\eta},\\
\end{cases}
\end{eqnarray}
and 
\begin{eqnarray}
\log(C(\alpha))= \begin{cases}
\log(\frac{2}{\eta}A^{\frac{2}{\eta}}) & \quad \alpha\leq \frac{2}{\eta}, \\
\log(\frac{2}{\eta}A^{\frac{2}{\eta}})+(\alpha-\frac{2}{2})\log(M)& \quad \alpha>\frac{2}{\eta}\\
\end{cases}
\end{eqnarray}
In other words, $\alpha=\frac{2}{\eta}$ is a first-order phase transition point
\begin{eqnarray}
\frac{\partial\log(C(\alpha))}{\partial \alpha}= \begin{cases}
0 & \quad \alpha\leq \frac{2}{\eta}, \\
\log(M)& \quad \alpha>\frac{2}{\eta}.\\
\end{cases}
\end{eqnarray}
There are several observations. First, according to Eq. \ref{pt1}, $M=\mathcal{O}(\text{poly}(n))$ can always enable efficient classical simulation. The slope and intercept of $\log(C(\alpha))$ tell us the information of $A$ and $M$.


\subsection{Beyond ideal case}
When $c_x^2$ (also $f(x)$) is not an ideal power function, we define the optimal matched power
%%%%%%%%%%%%%%%%%%%%%%%%%%%%%%%%%%%%%%%%%%%%%%%%%%%%%%%%%





\section{Relating error to Heisenberg Pauli back-propagated operator by random input states}

\section{Inverse-polynomial bounded spectra and complexity phase transition}

\section{From truncation to sampling}

\section{\szx{1}Heisenberg Pauli truncation}
Given a $n$-qubit noisy quantum circuit manifested as a CPTP map: $\mathcal{C}[\cdot]$ acting on an initial state $|\psi\rangle$, the goal is to estimate the expectation value of an operator $O$
\begin{eqnarray}
\langle O\rangle=\Tr(\mathcal{C}[|\psi\rangle\langle \psi|]O)=\Tr(|\psi\rangle\langle \psi|\mathcal{C}^\dag[O]).
\end{eqnarray}
We define $O_H=\mathcal{C}^\dag[O]$ and decompose it in the Pauli basis
\begin{eqnarray}
O_H=\mathcal{C}^\dag[O]=\sum_{i=1}^M c_i P_i,
\end{eqnarray}
where $\{c_i\}$ are arranged in descending order in terms of the absolute values. Under the vectorization picture, we have
\begin{eqnarray}\label{b3}
\langle O\rangle=(\langle \psi|\otimes \langle \psi^*|)\ket{O_H\rangle},
\end{eqnarray}
with $\ket{O_H\rangle}=\mathcal{V}[O_H]$ normalized.

Since $M$ can generally be exponentially large, the idea of truncation-based methods is to use another operator $\widetilde{O_H}$ of the form
\begin{eqnarray}
\widetilde{O_H}=\mathcal{C}^\dag[O]=\sum_{j=1}^{K-1} b_j P_j,
\end{eqnarray}
to approximate $O_H$. We also define normalized $\ket{\widetilde{O_H}\rangle}=\mathcal{V}[\widetilde{O_H}]$.

Consider an input state ensemble $\mu$ with $\|\mathbb{E}_{\psi\in\mu}|\psi\rangle\langle \psi|\|_{\infty}=2^{-n}g$ and define $\Delta O_{H}=O_H-\widetilde{O_H}$ and $\ket{\Delta O_{H}\rangle}=\ket{O_H\rangle}-\ket{\widetilde{O_H}\rangle}$, we have the average approximation error
\begin{eqnarray}\label{hpt}
\mathbb{E}_{\psi\in\mu}\Tr(|\psi\rangle\langle\psi| \Delta O_H)^2&&=\mathbb{E}_{\psi\in\mu}(\langle\psi|\otimes \langle\psi^*| )\ket{\Delta O_{H}\rangle}\bra{\langle\Delta O_{H}}(|\psi\rangle\otimes |\psi^*\rangle)\nonumber\\\leq && \mathbb{E}_{\psi\in\mu}\Tr((|\psi\rangle\langle \psi| \otimes I_n)\ket{\Delta O_{H}\rangle}\bra{\langle\Delta O_{H}})\nonumber\\\leq && \|\mathbb{E}_{\psi\in\mu}|\psi\rangle\langle \psi|\|_{\infty}\braket{\langle\Delta O_{H}|\Delta O_{H}\rangle}\nonumber\\=&&g \|\ket{\Delta O_{H}\rangle}\|_{n2}^2,
\end{eqnarray}
where the normalized vector-2 norm is defined as $\|\cdot\|_{n2}=2^{-n/2}\|\cdot\|_{2}$. For example, $\|\ket{O_H\rangle}\|_{n2}=\sum_i c_i^2$. When $\mu$ is a state 1-design, we have $g=1$ and the average approximation error is directly related to $\|\ket{\Delta O_{H}\rangle}\|_{n2}^2$.

\section{Pauli spectra phase transition diagnoses classical simulation ability}
Keeping $K-1\approx K$ Pauli terms in $\overline{O_H}$, the best we can expect or the smallest approximation error that can be achieved is to make all $b_i$ equal to $c_i$ resulting in
\begin{eqnarray}
\|\ket{\Delta O_{H}\rangle}\|_{n2}^2=\sum_{i=K}^M c_i^2=f(K).
\end{eqnarray}

Ignoring other pre/post-precessing procedures, we can understand $K$ as the complexity of Pauli truncation-based classical simulation methods. In this sense, the efficient classical simulation means that to achieve a $\epsilon$ error on average approximation \begin{eqnarray}
\mathbb{E}_{\psi\in\mu}\Tr(|\psi\rangle\langle\psi| \Delta O_H)^2=\epsilon^2,
\end{eqnarray}
we require 
\begin{eqnarray}
K=\mathcal{O}(\text{poly}(\epsilon^{-1}))\leq A \epsilon^{-\eta},
\end{eqnarray}
combining Eq. \ref{hpt}, the requirement needs
\begin{eqnarray}
f(K)\leq A^{\frac{2}{\eta}}K^{-\frac{2}{\eta}}.
\end{eqnarray}
In other words, efficient classical simulation indicates that $f(K)$ is bounded by a negative power function with a constant $A$.

\subsection{Ideal case}
Here, we assume $f(x)$ is exactly equal to $f(x)=A^{\frac{2}{\eta}}x^{-\frac{2}{\eta}}$, then we have
\begin{eqnarray}
-\frac{df(x)}{dx}*1=c_x^2=\frac{2}{\eta}A^{\frac{2}{\eta}}x^{-(1+\frac{2}{\eta})}=Bx^{-(1+\frac{2}{\eta})}=g(x),
\end{eqnarray}
with $B=\frac{2}{\eta}A^{\frac{2}{\eta}}$. Therefore, efficient classical simulation also indicates the Pauli spectra $\{c_i^2\}$ is bounded by a negative smaller than -1 power function with a constant $B$.

The value of $\eta$ can be diagnosed by numerical tests. Here, we show the strategies. We can try to use another function $h(\alpha,x)$ to bound $g(x)$
\begin{eqnarray}
h(\alpha, x)=C(\alpha)x^{-(1+\alpha)}\geq g(x),
\end{eqnarray}
leading to
\begin{eqnarray}
C(\alpha)=\max_x B x^{\alpha-\frac{2}{\eta}}=\max_x\frac{2}{\eta}A^{\frac{2}{\eta}}x^{\alpha-\frac{2}{\eta}}.
\end{eqnarray}
For different regimes, we have
\begin{eqnarray}\label{pt1}
C(\alpha)= \begin{cases}
\frac{2}{\eta}A^{\frac{2}{\eta}}x^{\alpha-\frac{2}{\eta}}\mid_{x=1}=\frac{2}{\eta}A^{\frac{2}{\eta}} & \quad \alpha\leq \frac{2}{\eta}, \\
\frac{2}{\eta}A^{\frac{2}{\eta}}x^{\alpha-\frac{2}{\eta}}\mid_{x=M}=\frac{2}{\eta}A^{\frac{2}{\eta}}M^{\alpha-\frac{2}{\eta}}& \quad \alpha>\frac{2}{\eta},\\
\end{cases}
\end{eqnarray}
and 
\begin{eqnarray}
\log(C(\alpha))= \begin{cases}
\log(\frac{2}{\eta}A^{\frac{2}{\eta}}) & \quad \alpha\leq \frac{2}{\eta}, \\
\log(\frac{2}{\eta}A^{\frac{2}{\eta}})+(\alpha-\frac{2}{\eta})\log(M)& \quad \alpha>\frac{2}{\eta}\\
\end{cases}
\end{eqnarray}
In other words, $\alpha=\frac{2}{\eta}$ is a first-order phase transition point
\begin{eqnarray}
\frac{\partial\log(C(\alpha))}{\partial \alpha}= \begin{cases}
0 & \quad \alpha\leq \frac{2}{\eta}, \\
\log(M)& \quad \alpha>\frac{2}{\eta}.\\
\end{cases}
\end{eqnarray}

There are several observations. First, according to Eq. \ref{pt1}, $M=\mathcal{O}(\text{poly}(n))$ can always enable efficient classical simulation. The slope and intercept of $\log(C(\alpha))$ tell us the information of $A$ and $M$.

\subsection{Non-ideal case}
When $g(x)$ (also $f(x)$) is not an ideal power function, we define $\eta_o$ as the optimal matched power
\begin{definition}[Optimal matched power]
For each $\eta>0$, we define 
\begin{eqnarray}
A_u(\eta)=\rm{arg}\min_{A} \frac{2}{\eta}A^{\frac{2}{\eta}}x^{-(1+\frac{2}{\eta})}\geq c_x^2,\forall x
\end{eqnarray}
\begin{eqnarray}
A_l(\eta)=\rm{arg}\max_{A} \frac{2}{\eta}A^{\frac{2}{\eta}}x^{-(1+\frac{2}{\eta})}\leq c_x^2,\forall x
\end{eqnarray}
then, the optimal matched power is defined as
\begin{eqnarray}
\eta_o=\rm{arg}\min_{\eta} A_u(\eta)/A_l(\eta).
\end{eqnarray}
\end{definition}

For non-ideal case, when we use $h(\alpha,x)$ to bound $g(x)$, we have
\begin{eqnarray}
h(\alpha, x)=C(\alpha)x^{-(1+\alpha)}\geq g(x)\geq \frac{2}{\eta_o}A_l(\eta_o)^{\frac{2}{\eta_o}}x^{-(1+\frac{2}{\eta_o})},\forall x
\end{eqnarray}
leading to
\begin{eqnarray}
C(\alpha)\geq\max_x  \frac{2}{\eta_o}A_l(\eta_o)^{\frac{2}{\eta_o}} x^{\alpha-\frac{2}{\eta_o}}.
\end{eqnarray}
Also, we have 
\begin{eqnarray}
h(\alpha, x)=C(\alpha)x^{-(1+\alpha)}\leq \frac{2}{\eta_o}A_u(\eta_o)^{\frac{2}{\eta_o}}x^{-(1+\frac{2}{\eta_o})},\exists x,
\end{eqnarray}
leading to 
\begin{eqnarray}
C(\alpha)\leq\max_x  \frac{2}{\eta_o}A_u(\eta_o)^{\frac{2}{\eta_o}} x^{\alpha-\frac{2}{\eta_o}}.
\end{eqnarray}

Therefore, we have
\begin{eqnarray}
\log(C(\alpha))\in \begin{cases}
[\log(\frac{2}{\eta_o}A_l(\eta_o)^{\frac{2}{\eta_o}},\log(\frac{2}{\eta_o}A_u(\eta_o)^{\frac{2}{\eta_o}}]& \quad \alpha\leq \frac{2}{\eta_o}, \\
[\log(\frac{2}{\eta_o}A_l(\eta_o)^{\frac{2}{\eta_o}}+(\alpha-\frac{2}{\eta_o})\log(M),\log(\frac{2}{\eta_o}A_u(\eta_o)^{\frac{2}{\eta_o}}+(\alpha-\frac{2}{\eta_o})\log(M)]& \quad \alpha>\frac{2}{\eta_o}\\
\end{cases}
\end{eqnarray}
The phase transition point now has an uncertainty $\frac{2}{\eta_o}\pm \delta $ with
\begin{eqnarray}
\delta =\frac{\log(\frac{2}{\eta_o}A_u(\eta_o)^{\frac{2}{\eta_o}}-\log(\frac{2}{\eta_o}A_l(\eta_o)^{\frac{2}{\eta_o}}}{\log(M)}=\frac{2\log(A_u(\eta_o)/A_l(\eta_o))}{\eta_o\log(M)}.
\end{eqnarray}
Therefore, optimal matched power $\eta_o$ has the smallest uncertainty, which makes sense of its definition.




\end{appendix}

\section{numerical simulation, case by case}

\end{document}
